\section{Lecture 9 (February 3, 2025, est.): Affine groups and tensor automorphisms}\label{lec9}
\noindent\textit{Scribe: Guillermo Sanmarco}


Our first reconstruction problem is to recover an affine group scheme $G$ from its finite-dimensional representation category together with the forgetful functor. Throughout this lecture $k$ is a field; no finite-type assumption on $G$ is imposed.

\subsection*{Approximation by algebraic groups}
\begin{proposition}\label{prop:hopf-finite-generation}
Every finite subset of a commutative Hopf algebra $A$ over $k$ is contained in a Hopf subalgebra finitely generated as a $k$-algebra.
\end{proposition}
\begin{proof}
By local finiteness of the regular comodule, the given subset lies in a finite-dimensional subcomodule $V\subseteq A$. Write $\Delta(v_j)=\sum_i v_i\otimes a_{ij}$ in a basis of $V$. Coassociativity gives
\[
 \Delta(a_{ij})=\sum_s a_{is}\otimes a_{sj}.
\]
Thus the finite-dimensional span $C$ of the $v_j$, the $a_{ij}$, and $1$ satisfies $\Delta(C)\subseteq C\otimes C$. Let $A_0$ be the algebra generated by $C+S(C)$. Since the antipode reverses coproducts, both $\Delta(C)$ and $\Delta(S(C))$ lie in $A_0\otimes A_0$. For a commutative Hopf algebra $S^2=\mathrm{id}$, so $S(A_0)\subseteq A_0$. The counit restricts, and $A_0$ is the desired Hopf subalgebra.
\end{proof}
The commutativity assumption matters for this short argument: it is not enough, for a completely arbitrary Hopf algebra, merely to add one image under the antipode.

\begin{corollary}\label{cor:group-inverse-limit}
Every affine group scheme is an inverse limit of affine algebraic group schemes
\[
 G\simeq\varprojlim_i G_i
\]
with faithfully flat transition morphisms.
\end{corollary}
\begin{proof}
Write $A=\mathcal O(G)=\bigcup_iA_i$, directed by inclusion over its finitely generated Hopf subalgebras. Then $G_i=\operatorname{Spec}A_i$ is algebraic, and
\[
 \operatorname{Hom}_{k\text{-alg}}(A,R)
 \simeq\varprojlim_i\operatorname{Hom}_{k\text{-alg}}(A_i,R)
\]
for every commutative $k$-algebra $R$. An inclusion of commutative Hopf algebras over a field is faithfully flat; applying this to $A_i\subseteq A_j$ gives the final assertion. This faithful-flatness theorem is proved in \cite[Chapter VI, Section 11]{milneAGS}.
\end{proof}
The algebras $A_i$ are finitely generated, not necessarily finite dimensional. Also, faithful flatness of $G_j\to G_i$ does not assert surjectivity of $G_j(R)\to G_i(R)$ for every test algebra $R$.

\subsection*{The tensor automorphism functor}
Let $\omega:\operatorname{Rep}_k(G)\to\mathrm{Vec}_k$ forget the action. For a commutative $k$-algebra $R$, put $\omega_R(V)=\omega(V)\otimes_kR$.
\begin{definition}
The group $\underline{\operatorname{Aut}}^{\otimes}(\omega)(R)$ consists of families of invertible $R$-linear maps $\lambda_V:\omega_R(V)\to\omega_R(V)$ satisfying
\begin{enumerate}
\item $(f\otimes\mathrm{id}_R)\lambda_V=\lambda_W(f\otimes\mathrm{id}_R)$ for every morphism $f:V\to W$;
\item $\lambda_{V\otimes W}=\lambda_V\otimes_R\lambda_W$ under the canonical tensor identification;
\item $\lambda_k=\mathrm{id}_R$ on the trivial representation.
\end{enumerate}
Composition is componentwise; scalar extension makes this a group-valued functor of $R$.
\end{definition}
An element $g\in G(R)$ acts on every $V_R$, and equivariance of morphisms and the diagonal tensor action give a natural group homomorphism
\[
 G\longrightarrow\underline{\operatorname{Aut}}^{\otimes}(\omega).
\]
We shall prove that it is an isomorphism in Lecture~\ref{lec10}.

\subsection*{Reduction to a tensor generator}
For $X\in\operatorname{Rep}_k(G)$, let $\mathcal C_X$ be the full subcategory of subquotients of finite direct sums of tensor products of $X$ and $X^*$, including the unit. A tensor automorphism of $\omega|_{\mathcal C_X}$ is determined by its component on $X$. Indeed, the component on $X^*$ is determined by evaluation, those on tensor products and sums follow from tensor compatibility and naturality, and those on subquotients follow from injectivity and surjectivity after tensoring with $R$. Hence
\[
 \underline{\operatorname{Aut}}^{\otimes}(\omega|_{\mathcal C_X})(R)
 \hookrightarrow\mathrm{GL}(X\otimes R).
\]

Let $G_X$ be the scheme-theoretic image of $G\to\mathrm{GL}(X)$. It is a closed algebraic subgroup. Its coordinate Hopf algebra is generated by the matrix coefficients of $X$ and its dual. A faithful finite-dimensional representation generates the representation category of an algebraic group under tensor products, sums, duals and subquotients; see \cite[Chapter VIII, Section 11]{milneAGS}. Thus $\mathcal C_X$ identifies with $\operatorname{Rep}(G_X)$, and its action gives inclusions
\[
 G_X(R)\subseteq
 \underline{\operatorname{Aut}}^{\otimes}(\omega|_{\mathcal C_X})(R)
 \subseteq\mathrm{GL}(X\otimes R).
\]

Naturality gives preservation of every subrepresentation and compatibility with every quotient. For example, if $Y\subseteq V$, the inclusion gives
\[
 \lambda_V(Y\otimes R)=Y\otimes R.
\]
If $t\in V^G$, the morphism $k\to V$, $1\mapsto t$, gives $\lambda_V(t\otimes1)=t\otimes1$.

Finally, as $X$ varies and direct sums provide upper bounds, the Hopf algebras $\mathcal O(G_X)$ exhaust $\mathcal O(G)$. Indeed, every element belongs to a finite-dimensional regular subcomodule and is a linear combination of its matrix coefficients by the counit identity. Thus
\[
 G=\varprojlim_XG_X,\qquad
 \underline{\operatorname{Aut}}^{\otimes}(\omega)
 =\varprojlim_X\underline{\operatorname{Aut}}^{\otimes}(\omega|_{\mathcal C_X}).
\]
The second equality holds because every finite collection of representations and morphisms lies in some $\mathcal C_X$.
